\begin{table}[htbp]
\centering
\small
\caption{S0--S5 虚拟台风情景事件级核心指标与相对变化}
\label{tab:problem3_scenario_comparison}
\resizebox{\textwidth}{!}{
\begin{tabular}{c l r r r r r r p{4.3cm}}
\hline
情景 & 情景名称 & 最大累计/mm & 最大$A_{10}$/km$^2$ & 最大$D_{10}$/h & $\Delta C_{\max}$/\% & $\Delta A_{10}$/\% & $\Delta D_{10}$/\% & 简要结论 \\
\hline
S0 & baseline\_kong\_rey & 723.48 & 74508.15 & 24.5 & 0.00 & 0.00 & 0.00 & 问题三统一链条下的基准对照。 \\
S1 & intensity\_enhanced & 746.93 & 79501.25 & 26.0 & 3.24 & 6.70 & 6.12 & 强度增强主要抬升强降水面积和雨强尾部。 \\
S2 & near\_coast\_westward\_path\_shift & 798.34 & 71587.80 & 26.5 & 10.35 & -3.92 & 8.16 & 路径近岸改变落区重叠，强降水面积不单调增加。 \\
S3 & nearshore\_segment\_landfall\_shift & 705.52 & 74772.72 & 23.5 & -2.48 & 0.36 & -4.08 & 近岸段偏移主要体现沿海影响带转移。 \\
S4 & slower\_translation & 1208.09 & 75003.97 & 39.0 & 66.98 & 0.67 & 59.18 & 移速减慢显著放大累计降水和固定格点持续时间。 \\
S5 & compound\_moderate\_high\_risk & 947.66 & 82748.91 & 31.5 & 30.99 & 11.06 & 28.57 & 复合扰动使多项风险指标同步抬升。 \\
\hline
\end{tabular}
}
\end{table}
